\section*{Chapter \ref{ch:mx:the-limit-order-book} --- The Limit Order Book}

\begin{solution}{exo:mx:the-limit-order-book:1}
The sell executes against the level in time order: 300 shares against the first order and 150 against the second. The best bid stays 100.00 with
$50+200=250$ shares in two orders; the best ask is unchanged at 100.01.
\end{solution}

\begin{solution}{exo:mx:the-limit-order-book:2}
$P_3=\tfrac{1}{1.2}\cdot\tfrac{1.1}{1.3}\cdot\tfrac{1.2}{1.4}\cdot\tfrac{1.3}{1.5}=0.524$.
\end{solution}

\begin{solution}{exo:mx:the-limit-order-book:3}
The reduction keeps the order's place (an X message of 200 shares on an ITCH-style feed); the increase changes its priority on CME Globex, which puts
it behind every order at its price. On Nasdaq the increase is a cancel-replace: an Order Replace message removes the old reference and adds a new
one at the back.
\end{solution}

\begin{solution}{exo:mx:the-limit-order-book:4}
From 99.00 to 101.00 in cents: $2\times100+1=201$ slots. A 3\% move takes the best price out of the band: the book must recentre (rebuild the array
around the new price) or keep the orders outside the band in an overflow map, and the recentring is a latency spike exactly when the market moves.
\end{solution}

\begin{solution}{exo:mx:the-limit-order-book:5}
$1/0.5+\sum_{k=1}^{10}1/(0.5+0.02k)=2+16.54=18.5$ seconds.
\end{solution}

\begin{solution}{exo:mx:the-limit-order-book:6}
The best level is where every execution happens and where stale orders are pulled first when the price moves through it; one tick behind, orders
accumulate out of reach of small \omterm{def:mx:the-limit-order-book:orders}{market orders} and are replenished by the orders that arrive there, so they pile up. The shape is the hump
chapter~3 finds in real books.
\end{solution}

\begin{solution}{exo:mx:the-limit-order-book:7}
With \omterm{def:mx:the-limit-order-book:orders}{market orders} at rate 1 the spread falls from 2.55 to 1.64 ticks and the best level holds 3.29 orders instead of 2.72: \omterm{def:mx:the-limit-order-book:orders}{market orders} empty the
best level half as often. Behind five orders the fill probability is 0.602 against the proposition's 0.769, a loss of 22\% (23\% at rate 2): waits
are twice as long, which gives later orders more time to improve the price, but a narrower spread leaves fewer prices to improve into, and the two
effects nearly offset.
\end{solution}

\begin{solution}{exo:mx:the-limit-order-book:8}
The factor $\mu/(\mu+\nu)$ assumes nobody can take the order's place at the front. Any bid placed above it becomes the best and takes the market
sells first; the level may stop being the best, and the order then waits indefinitely. Ten seconds at the front also says something about the
rates: in a quiet market $\mu$ is lower than its average.
\end{solution}

\begin{solution}{pb:mx:the-limit-order-book:1}
\textbf{1.} 31\,535 messages: 49.7\% adds, 47.8\% cancellations, 2.5\% executions.
\textbf{2.} 40 messages per trade; 3.8\% of the orders that left the book did so by trading.
\textbf{3.} 15.5 seconds in the median; 24.9 for orders that traded, 15.3 for those cancelled: trading takes longer than giving up.
\textbf{4.} 4.5 seconds for the best bid (4.8 for the best ask).
\textbf{5.} 2\,834 orders, with a median of 1\,400 shares ahead.
\textbf{6.} 13\% trade at least in part, 10\% in full.
\textbf{7.} One tick behind the best (1\,508 shares on the bid side against 1\,366 at the best): the best level is emptied by executions and stale
cancellations, the next one is not.
\textbf{8.} Order lives and the share of orders that trade would rise; queues ahead would be longer and depth larger; messages per trade would fall.
\textbf{9.} Unit orders; \omterm{def:mx:the-limit-order-book:orders}{market orders} at rate $\mu$ from the front; cancellations at $\theta$ per order ahead and $\nu$ for ours; independent
exponential clocks; nobody jumps ahead. Then $P_n=\tfrac{\mu}{\mu+\nu}\prod_{k=1}^{n}\tfrac{\mu+k\theta}{\mu+k\theta+\nu}$.
\textbf{10.} 0.952, 0.870 and 0.755.
\textbf{11.} \texttt{queue\_mc(5, 2, 0.05, 0.05)} with 20\,000 paths agrees to within 0.01 (0.869).
\textbf{12.} 2.83 seconds.
\textbf{13.} 0.953 behind one, 0.871 behind five: the best bid queue then behaves exactly as the model assumes, because nobody can place a bid above
it.
\textbf{14.} 0.912, 0.672 and 0.535.
\textbf{15.} \emph{Named result}: the fill probability behind $n$ orders falls from 0.95 ($n=1$) to 0.76 ($n=12$) in the proposition; in the
zero-intelligence market with price improvement it falls from 0.91 to 0.54 behind eight, and improvement takes away 23\% of the promised fills behind
five and 35\% behind eight.
\textbf{16.} Only \omterm{def:mx:the-limit-order-book:orders}{market orders} and cancellations remove orders from the best levels, and nothing is ever placed inside the spread, so each side's
best price only moves away.
\textbf{17.} 2.55 ticks and 2.72 orders.
\textbf{18.} That nobody gets ahead (price improvement), together with constant rates: both bias the proposition upwards; in a real book, flow at the
best also clusters in time (One Quant Book 4, chapter~7).
\textbf{19.} $\mu$ from the executions at the best per unit time, $\theta$ from cancellations at the best divided by the time-integrated queue size,
$\nu$ from one's own order lifetimes; all conditional on the queue's state.
\textbf{20.} The probability, and the time, of filling before the price moves: the front of a long queue is worth a fill; the back of it, a lottery.
\end{solution}

\begin{solution}{iq:mx:the-limit-order-book:1}
A hash map from order reference to order (cancels and executions name it), a sorted map or price-indexed array from price to level, and a FIFO
list per level. Cancel and execute: $O(1)$ with an intrusive list; add: $O(1)$, or $O(\log L)$ when a tree gets a new level; best price: $O(1)$ with a
cached pointer. \iqlookfor{the three access paths, the hash map for references, why arrays win when prices cluster.}
\end{solution}

\begin{solution}{iq:mx:the-limit-order-book:2}
A decrease only removes the order's own shares, which harms nobody behind it; an increase would let a trader hold the front with a small order and
enlarge it when a big counterparty appears, jumping everyone who waited. \iqlookfor{fairness to the orders behind, and the gaming a size increase
would allow.}
\end{solution}

\begin{solution}{iq:mx:the-limit-order-book:3}
With four ahead: $P=\tfrac{\mu}{\mu+\nu}\prod_{k=1}^{4}\tfrac{\mu+k\theta}{\mu+k\theta+\nu}$, from competing exponential clocks and memorylessness.
\iqlookfor{the stage decomposition and the independence of stages; a remark that price improvement lowers it.}
\end{solution}

\begin{solution}{iq:mx:the-limit-order-book:4}
When the price leaves the band the array covers (a fast market, a halt reopening far away), or for instruments with huge price ranges and sparse
books. Recentre the array, keep an overflow map for far orders, and size the band from the instrument's volatility. \iqlookfor{memory against range,
and the latency spike of recentring at the worst moment.}
\end{solution}

\begin{solution}{iq:mx:the-limit-order-book:5}
Not in itself: cancellations are how liquidity providers update quotes when information arrives, and the displayed book stays current. It costs
message capacity and makes displayed depth less firm (chapter~18's quote fade); regulators watch it through order-to-trade ratios (chapter~24).
\iqlookfor{a two-sided answer: stale quotes are worse, but fleeting liquidity misleads.}
\end{solution}

\begin{solution}{iq:mx:the-limit-order-book:6}
Right: the order of magnitude of the spread and its scaling with the order-flow rates, the shape of depth near the best, why most orders do not trade.
Wrong: order flow is not independent (signs of trades have long memory, chapter~12), rates react to the state and to information, and prices have no
anchor to value. \iqlookfor{it is a null model: a benchmark for what needs no intelligence to explain.}
\end{solution}
